\section{Related Literature}
\label{sec:literature}

This paper contributes to several strands of the economics literature.

\paragraph{Government spending efficiency and procurement.}
A foundational distinction in the study of public spending separates \textit{active waste} (corruption) from \textit{passive waste} (inefficiency), as formalized by \citet{bandiera2009active}, who estimate that passive waste accounts for the majority of excess spending in Italian public procurement. Subsequent work has identified specific channels of passive waste, including bureaucratic capacity \citep{best2023procurement}, management practices \citep{rasul2018management, bloom2015impact}, and the design of procurement mechanisms \citep{lewisfaupel2016can, szucs2023competition}. Our paper identifies a previously unstudied source of passive waste: the compression of planning time induced by judicial mandates. We show that this channel operates through both the reference price setting process (the internal phase) and the competitive bidding process (the external phase).

\paragraph{Procurement design and competition.}
A growing literature documents how procurement design affects competition and prices. \citet{coviello2018bidding} show that court-imposed delays reduce bidder participation and increase procurement costs in Italian public works. \citet{decarolis2018buyer} demonstrate that buyer competence is a key determinant of procurement outcomes. \citet{baltrunaite2021discretion} find that reducing discretion in procurement can improve outcomes when corruption is prevalent but may harm outcomes otherwise. Our setting differs in that judicial intervention compresses rather than extends the procurement timeline, and the primary mechanism operates through planning disruption rather than direct manipulation of award criteria.

\paragraph{Accountability, oversight, and unintended consequences.}
A related literature examines how accountability mechanisms intended to reduce waste can themselves generate costs. \citet{rasul2018management} show that management practices in the Nigerian civil service affect project completion rates, while \citet{williams2021bureaucratic} finds that management quality explains substantial variation in public service delivery. \citet{duflo2018obligations} argue that the design of incentives for public officials must balance accountability with operational flexibility. Our ``under the gun'' effect is a specific instance of this tension: judicial sanctions intended to ensure compliance with health mandates inadvertently raise procurement costs by distorting the bargaining environment.

\paragraph{Health litigation and right-to-health enforcement.}
The legal and public health literatures have documented the growth of health litigation in developing countries, particularly in Latin America \citep{wang2015right, ferraz2009right, biehl2009will}. \citet{cnj2019judicializacao} provide comprehensive data on the judicialization of health in Brazil. However, the existing literature has focused primarily on the legal and public health dimensions, with limited attention to the fiscal costs transmitted through the procurement channel. Our paper fills this gap by providing granular estimates of how judicial enforcement affects specific procurement outcomes---prices, quantities, competition, and tender success.

\paragraph{Auction theory under urgency.}
From a theoretical perspective, our findings connect to the auction theory literature on how the number of bidders and time pressure affect prices. \citet{bulow1996auctions} show that attracting an additional bidder can be more valuable than optimizing the auction mechanism. When judicial mandates reduce bidder participation, they directly undermine the competitive benefits of procurement auctions. Moreover, the public information that a tender is court-mandated creates an asymmetric bargaining environment that further erodes the government's position \citep{krishna2009auction}.
